\subsection{Roots of a polynomial in one indeterminate. Multiplicity}

Let \(g \in A[(X_i)_{i \in I}]\) and let \(E\) be a unital associative \(A\) -algebra. Let \(x = (x_i)_{i \in I}\) be a family of pairwise permutable elements of \(E\) . We shall say that \(x\) is a zero of \(g\) in \(E^{I}\) if \(g(x) = 0\) . If \(f\) is a polynomial in a single indeterminate, a zero of \(f\) in \(E\) is also called a root of \(f\) in \(E\) .

PROPOSITION 1. --- Let \(f \in \mathbf{A}[\mathbf{X}]\) and \(\alpha \in \mathbf{A}\) . The remainder of the division of \(f\) by \(\mathbf{X} - \alpha\) is \(f(\alpha)\) . For \(\alpha\) to be a root of \(f\) it is necessary and sufficient that \(\mathbf{X} - \alpha\) should be a divisor of \(f\) in \(\mathbf{A}[\mathbf{X}]\) .

For if \(u, v \in \mathbf{A}[\mathbf{X}]\) are such that \(f = (\mathbf{X} - \alpha)u + v\) , \(\deg v < 1\) , then \(v\) is a scalar and \(f(\alpha) = (\alpha - \alpha)u(\alpha) + v = v\) . This proves the first assertion and the second follows from it.

PROPOSITION 2. --- Let \(f \in \mathbf{A}[\mathbf{X}]\) , \(\alpha \in \mathbf{A}\) , and let \(h\) be an integer \(\geqslant 0\) . The following conditions are equivalent:

\begin{enumerate}
\def\labelenumi{(\roman{enumi})}
\item
  \(f\) is divisible by \((\mathbf{X} - \alpha)^h\) but not by \((\mathbf{X} - \alpha)^{h + 1}\) ;
\item
  there exists \(g \in \mathbf{A}[\mathbf{X}]\) such that \(f = (\mathbf{X} - \alpha)^h g\) and \(g(\alpha) \neq 0\) .
\item
  \(\Rightarrow\) (ii) follows at once from Prop. 1.
\item
  \(\Rightarrow\) (i): Suppose that \(f = (\mathbf{X} - \alpha)^h g\) , where \(g\) does not admit \(\alpha\) as a root. Then \(f\) is divisible by \((\mathbf{X} - \alpha)^h\) ; if \(g_1 \in \mathbf{A}[\mathbf{X}]\) existed such that \(f = (\mathbf{X} - \alpha)^{h + 1}g_1\) , then since \((\mathbf{X} - \alpha)^n\) is not a divisor of zero in \(\mathbf{A}[\mathbf{X}]\) (IV, p.~9, Prop. 7), we have \(g = (\mathbf{X} - \alpha)g_1\) and so \(g(\alpha) = 0\) , which is absurd.
\end{enumerate}

PROPOSITION 3. --- Let \(f\) be a non-zero element of \(\mathbf{A}[\mathbf{X}]\) and \(\alpha \in \mathbf{A}\) . There exists just one integer \(h \geqslant 0\) satisfying the conditions (i) and (ii) of Prop. 2.

This is clear for condition (i), bearing in mind the fact that if \(f\) is divisible by \((\mathbf{X} - \alpha)^h\) , then \(\deg f \geqslant h\) (IV, p.~9, Prop. 7).

DEFINITION 1. --- With the above notation we say that \(\alpha\) is of order \(h\) , or multiplicity \(h\) relative to \(f\) .

If h \textgreater{} 0 we also say that \(\alpha\) is a root of order h or multiplicity h of f.~A root of order 1 is called a simple root, a root of order 2 a double root,\ldots{} A root of order \textgreater{} 1 is said to be multiple.

Remarks. --- 1) If \(f=0\) we agree to say that \(\alpha\) has order \(\geqslant h\) relative to \(f\) , whatever \(\alpha\in\mathbf{A}\) and the integer \(h\geqslant0\) . For any \(f\in\mathbf{A}[\mathbf{X}]\) and \(\alpha\in\mathbf{A}\) , to say that \(\alpha\) has order \(\geqslant h\) relative to \(f\) means that \((\mathbf{X}-\alpha)^{h}\) divides \(f\) .

\begin{enumerate}
\def\labelenumi{\arabic{enumi})}
\setcounter{enumi}{1}
\tightlist
\item
  Let B be a commutative ring containing A as subring. Let \(f \in A[X]\) be nonzero and \(\alpha \in A\) . The order of \(\alpha\) relative to f is the same, whether we consider f as element of B{[}X{]} or as element of A{[}X{]}. This is clear from condition (ii) of Prop. 2.
\end{enumerate}

PROPOSITION 4. --- Let f and g be non-zero elements of A{[}X{]}. Let \(\alpha \in \mathbf{A}\) , and let the orders of \(\alpha\) relative to f and g be p and q respectively.

\begin{enumerate}
\def\labelenumi{(\roman{enumi})}
\item
  The order of \(\alpha\) relative to \(f + g\) is \(\geqslant \inf(p, q)\) . It is equal to \(\inf(p, q)\) if \(p \neq q\) .
\item
  The order of \(\alpha\) relative to \(fg\) is \(\geqslant p + q\) . It is equal to \(p + q\) if \(\mathbf{A}\) is an integral domain.
\end{enumerate}

For we have \(f(\mathbf{X}) = (\mathbf{X} - \alpha)^p f_1(\mathbf{X}), g(\mathbf{X}) = (\mathbf{X} - \alpha)^q g_1(\mathbf{X})\) with \(f_1(\alpha) \neq 0\) , \(g_1(\alpha) \neq 0\) . Suppose for example that \(p \leqslant q\) ; then we have

\[
f (\mathbf {X}) + g (\mathbf {X}) = (\mathbf {X} - \alpha) ^ {p} \left(f _ {1} (\mathbf {X}) + (\mathbf {X} - \alpha) ^ {q - p} g _ {1} (\mathbf {X})\right),
\]

and if \(p < q\) , \(\alpha\) is not a root of \(f_1(\mathbf{X}) + (\mathbf{X} - \alpha)^{q - p} g_1(\mathbf{X})\) ; this proves (i). On the other hand, we have \(f(\mathbf{X})g(\mathbf{X}) = (\mathbf{X} - \alpha)^{p + q} f_1(\mathbf{X})g_1(\mathbf{X})\) and \(f_1(\alpha)g_1(\alpha) \neq 0\) if \(\mathbf{A}\) is an integral domain; this proves (ii).

PROPOSITION 5. --- Suppose that A is an integral domain. Let f be a non-zero element of A{[}X{]}, and \(\alpha_{1},\ldots,\alpha_{p}\) pairwise distinct roots of f in A, of orders \(k_{1},\ldots,k_{p}\) . We have

\[
f (\mathrm{X}) = \left(\mathrm{X} - \alpha_ {1}\right) ^ {k _ {1}} \left(\mathrm{X} - \alpha_ {2}\right) ^ {k _ {2}} \dots \left(\mathrm{X} - \alpha_ {p}\right) ^ {k _ {p}} g (\mathrm{X})
\]

where \(g \in \mathbf{A}[\mathbf{X}]\) and \(\alpha_1, \ldots, \alpha_p\) are not roots of \(g\) .

We proceed by induction on \(p\) , the proposition being evident for \(p = 1\) , by Def. 1. Suppose then that \(f(\mathbf{X}) = g_1(\mathbf{X})g_2(\mathbf{X})\) , where

\[
g _ {1} (\mathrm{X}) = \left(\mathrm{X} - \alpha_ {1}\right) ^ {k _ {1}} \dots \left(\mathrm{X} - \alpha_ {p - 1}\right) ^ {k _ {p - 1}}, \quad g _ {2} (\mathrm{X}) \in \mathrm{A} [ \mathrm{X} ].
\]

Since A is an integral domain and \(\alpha_{p}\) is distinct from \(\alpha_{1},\ldots,\alpha_{p-1}\) it follows that \(\alpha_{p}\) is not a root of \(g_{1}(X)\) , hence \(\alpha_{p}\) is a root of order \(k_{p}\) of \(g_{2}(X)\) (Prop. 4, (ii)). It follows that \(g_{2}(X)\) is divisible by \((\mathbf{X}-\alpha_{p})^{k_{p}}\) , and so

\[
f (\mathbf {X}) = \left(\mathbf {X} - \alpha_ {1}\right) ^ {k _ {1}} \dots \left(\mathbf {X} - \alpha_ {p}\right) ^ {k _ {p}} g (\mathbf {X})
\]

where \(g(\mathbf{X}) \in \mathbf{A}[\mathbf{X}]\) . Clearly \(\alpha_1, \ldots, \alpha_p\) are not roots of \(g\) .

THEOREM 1. --- Let A be an integral domain. Given a non-zero element f of A{[}X{]}, of degree n, then the sum of the orders of all the roots of f in A is ≤ n. This follows immediately from Prop. 5.

COROLLARY. --- Assume that A is an integral domain and let \(f, g \in A[X]\) , of degrees \(\leqslant n\) . If there exist \(n+1\) pairwise distinct elements \(x_{1}, \ldots, x_{n+1}\) of A such that \(f(x_{i}) = g(x_{i})\) for \(1 \leqslant i \leqslant n+1\) , then \(f = g\) .

It suffices to apply Th. 1 to f - g.

PROPOSITION 6 (Lagrange interpolation formula). --- Let K be a commutative field, \(\alpha_{1}, \alpha_{2}, \ldots, \alpha_{n}\) distinct elements of K and \(\beta_{1}, \beta_{2}, \ldots, \beta_{n}\) any elements of K. For \(i = 1, 2, \ldots, n\) we put

\[
f _ {i} (\mathbf {X}) = \prod_ {j \in U (i)} \left(\mathbf {X} - \alpha_ {j}\right) / \left(\alpha_ {i} - \alpha_ {j}\right),
\]

where \(\mathrm{U}(i)\) is the set of integers j such that \(j \neq i\) and \(1 \leqslant j \leqslant n\) . Then \(\beta_{1}f_{1} + \cdots + \beta_{n}f_{n}\) is the unique element f of K{[}X{]} such that \(\deg f < n\) and \(f(\alpha_{i}) = \beta_{i}\) for \(1 \leqslant i \leqslant n\) .

The uniqueness of \(f\) follows from the Cor. to Th. 1. Let \(f = \beta_{1}f_{1} + \cdots + \beta_{n}f_{n}\) , then since \(f_{i}\) has degree \(n - 1\) , we have \(\deg f < n\) . On the other hand, \(f_{i}(\alpha_{j}) = 0\) for \(j \neq i\) and \(f_{i}(\alpha_{i}) = 1\) , hence \(f(\alpha_{i}) = \beta_{i}\) for \(1 \leqslant i \leqslant n\) .

COROLLARY. --- Suppose that A is an integral domain. Let \(f \in A[X]\) , of degree \textless{} n, and let K be a subring of A which is a field. If there exist n distinct elements \(\alpha_{1}, \ldots, \alpha_{n}\) of A such that \(\alpha_{i} \in K\) and \(f(\alpha_{i}) \in K\) for \(i = 1, \ldots, n\) , then \(f \in K[X]\) .

